%============================================================
% CH1 · PATTERN 2 — EXERCISES, HINTS, ANSWERS
%============================================================
\subsection{Pattern 2: Exercises and Going Further}
\label{ssec:ch1-p2-ex}

\begin{patternbox}{Pattern 2 Summary Card}
\begin{tabularx}{\linewidth}{@{}>{\bfseries}l X@{}}
Question & which vertices are reachable from $s$? --- the reach set $R(s)$ \\
Algorithm & mark $s$; repeatedly take a vertex from a stack/queue and mark + add its unmarked out-neighbours; $\VE$ \\
Proof & marked $\Rightarrow$ reachable (induction); at the end no edge leaves the marked set $\Rightarrow$ all reachable are marked \\
Key variations & yes/no with early exit, ``all reached?'' (count $= n$), size of $R(s)$, many sources, reverse graph, one search per vertex \\
Top traps & undirected edge stored once; wrong edge direction; deep recursion; marking on pop; counting edges instead of vertices \\
\end{tabularx}
\end{patternbox}

\exlegend
\begin{enumerate}[leftmargin=*, label=\textbf{E2.\arabic*}, itemsep=6pt]
    \item \exeasy\ \pconcept \cf{1000}\ \exprob{New Year Transportation}{Codeforces 500A}{https://codeforces.com/problemset/problem/500/A}. Cells $1..n$ ($n \le 3\cdot10^4$); from cell $i < n$ a one-way portal leads to cell $i + a_i$ ($1 \le a_i \le n-i$). Starting in cell 1, can you reach cell $t$? Print \texttt{YES} or \texttt{NO}.
    \item \exeasy\ \ptricky \exprob{Minimum Number of Vertices to Reach All Nodes}{LeetCode 1557}{https://leetcode.com/problems/minimum-number-of-vertices-to-reach-all-nodes/}\ \lc{Medium}. Given a DAG with $n \le 10^5$ vertices and up to $10^5$ edges, return the smallest set of vertices from which every vertex is reachable (it is unique).
    \item \exmedium\ \pconcept \exprob{Tour}{AtCoder ABC204 C}{https://atcoder.jp/contests/abc204/tasks/abc204_c}. A directed graph with $N \le 2000$ cities and $M \le 2000$ one-way roads. Count the ordered pairs (origin, destination) such that the destination can be reached from the origin using zero or more roads.
    \item \exmedium\ \pclassic \exprob{Reorder Routes to Make All Paths Lead to the City Zero}{LeetCode 1466}{https://leetcode.com/problems/reorder-routes-to-make-all-paths-lead-to-the-city-zero/}\ \lc{Medium}. $n \le 5\cdot10^4$ cities joined by $n-1$ one-way roads whose undirected version is a tree. Reverse as few roads as possible so that every city can reach city 0; return that number.
    \item \exmedium\ \ptricky \exprob{Detonate the Maximum Bombs}{LeetCode 2101}{https://leetcode.com/problems/detonate-the-maximum-bombs/}\ \lc{Medium}. $n \le 100$ bombs $[x_i, y_i, r_i]$ ($1 \le x_i, y_i, r_i \le 10^5$). A detonated bomb detonates every bomb within Euclidean distance $r_i$ of its centre, and so on. Detonating exactly one bomb by hand, what is the maximum number of bombs that explode?
    \item \exmedium\ \pconcept \exown{Vertices on some $s$--$t$ walk.} Given a directed graph ($|V|, |E| \le 2\cdot 10^5$) and two vertices $s, t$, count the vertices $v$ such that some walk from $s$ to $t$ passes through $v$. On the running example $G$ of the theory section with $s = 5$, $t = 2$ the answer is $3$.
    \item \exmedium\ \precent \exprob{Properties Graph}{LeetCode 3493}{https://leetcode.com/problems/properties-graph/}\ \lc{Medium}. $n \le 100$ arrays $\texttt{properties}[i]$ of length $m \le 100$ with values in $1..100$, and an integer $k$. Vertices $i \ne j$ are joined by an undirected edge iff the two arrays have at least $k$ \emph{distinct} common values. Return the number of connected components.
\end{enumerate}

\subsubsection*{Hints}
\begin{enumerate}[leftmargin=*, label=\textbf{E2.\arabic*}]
    \item Every cell has exactly one outgoing portal, and portals only go to the right. What does a ``search'' look like when there is never a choice?
    \item In a DAG, which vertices can no other vertex reach? Must they be in the set? Is that already enough?
    \item $N + M$ is small. How expensive is one search? How many do you need?
    \item Search from city 0 over the roads \emph{ignoring} their direction. Every road you walk over is either already pointing towards 0 or not.
    \item The ``detonates'' relation is not symmetric: a big bomb can reach a small one but not the other way round. Build a directed graph, then search from every bomb.
    \item $v$ is on some $s$--$t$ walk iff $s \leadsto v$ and $v \leadsto t$. One search answers the first condition for all $v$; what answers the second?
    \item First build the edges (careful: ``distinct'' common values, so $[1,1]$ and $[1,1]$ share only one). Then: how many times does ``search from every vertex that is still unmarked'' start a new search?
\end{enumerate}

\subsubsection*{Answers}

\begin{solutionbox}[E2.1 New Year Transportation]
\textbf{Model:} vertex $=$ cell, edge $i \to i + a_i$; out-degree exactly 1 for $i < n$, so $R(1)$ is a single chain $1 \to 1 + a_1 \to \cdots$ (a \emph{functional graph}, Pattern~5). Positions strictly increase, so the walk ends after at most $n - 1$ jumps; $t$ is reachable iff the walk lands exactly on it.
\begin{lstlisting}[style=cpp]
#include <bits/stdc++.h>
using namespace std;

int main() {
    int n, t;
    scanf("%d %d", &n, &t);
    vector<int> a(n + 1, 0);
    for (int i = 1; i <= n - 1; i++) scanf("%d", &a[i]);
    int cur = 1;                          // every cell has exactly one portal: just follow it
    while (cur < t) cur += a[cur];        // positions only increase, so at most n-1 jumps
    puts(cur == t ? "YES" : "NO");
}
\end{lstlisting}
Samples (checked): $n=8, t=4$, $a = 1\,2\,1\,2\,1\,2\,1$ visits $1, 2, 4$ $\Rightarrow$ \texttt{YES}; $n=8, t=5$, $a = 1\,2\,1\,2\,1\,1\,1$ visits $1, 2, 4, 6$ and jumps over 5 $\Rightarrow$ \texttt{NO}. $\bigO{n}$.
\end{solutionbox}

\begin{solutionbox}[E2.2 Minimum Number of Vertices to Reach All Nodes]
A vertex with in-degree 0 can be reached only from itself, so it \emph{must} be in the set. Conversely, in a DAG every vertex $v$ with in-degree $> 0$ is reachable from some in-degree-0 vertex: walk backwards along incoming edges; without cycles the backward walk cannot repeat a vertex, so it stops --- at a vertex with in-degree 0. Hence the answer is exactly the set of in-degree-0 vertices: a reachability question solved by \textbf{Pattern~1}, without any search.
\begin{lstlisting}[style=cpp]
#include <bits/stdc++.h>
using namespace std;

class Solution {
public:
    vector<int> findSmallestSetOfVertices(int n, vector<vector<int>>& edges) {
        vector<char> hasIn(n, 0);
        for (auto& e : edges) hasIn[e[1]] = 1;
        vector<int> res;
        for (int v = 0; v < n; v++)
            if (!hasIn[v]) res.push_back(v);
        return res;
    }
};
\end{lstlisting}
Samples: $[0,3]$ and $[0,2,3]$ (checked). A brute force over all subsets agreed on 3000 random DAGs. In a graph \emph{with} cycles the rule breaks (a cycle has no in-degree-0 vertex); there one needs one vertex per source strongly connected component (Chapter~19).
\end{solutionbox}

\begin{solutionbox}[E2.3 Tour]
One search per origin: $N$ searches of $\bigO{N + M}$ each, i.e.\ $2000 \cdot 4000 = 8\cdot 10^6$ steps. The answer is $\sum_s |R(s)|$, and $(s,s)$ is counted because a journey may use zero roads.
\begin{lstlisting}[style=cpp]
#include <bits/stdc++.h>
using namespace std;

int main() {
    int n, m;
    scanf("%d %d", &n, &m);
    vector<vector<int>> adj(n);
    for (int i = 0; i < m; i++) {
        int a, b;
        scanf("%d %d", &a, &b);
        adj[a - 1].push_back(b - 1);                 // one-way road, 1-indexed input
    }
    long long answer = 0;
    vector<char> seen(n);
    vector<int> st;
    for (int s = 0; s < n; s++) {                    // one search per origin
        fill(seen.begin(), seen.end(), 0);
        seen[s] = 1; st.assign(1, s);
        int cnt = 1;                                 // (s, s) counts: zero roads
        while (!st.empty()) {
            int u = st.back(); st.pop_back();
            for (int v : adj[u])
                if (!seen[v]) { seen[v] = 1; cnt++; st.push_back(v); }
        }
        answer += cnt;
    }
    cout << answer << "\n";
}
\end{lstlisting}
Samples (checked): $7$, $3$, $16$. In Sample~1 (roads $1{\to}2$, $2{\to}3$, $3{\to}2$) the reach sets have sizes $3, 2, 2$. The theory section's \texttt{reachable pairs = 30} is the same computation on $G$.
\end{solutionbox}

\begin{solutionbox}[E2.4 Reorder Routes to Make All Paths Lead to the City Zero]
Store each road $a \to b$ twice: as $a \to b$ with cost 1 (it points \emph{away} from where we came from) and as $b \to a$ with cost 0. Search from city 0 over this undirected tree. When the search crosses a road from $u$ to a new city $v$, the road must end up pointing $v \to u$ (towards 0); it costs 1 exactly if it was stored as $u \to v$. Every road is crossed exactly once (the graph is a tree), so the sum of costs is the answer --- and no smaller answer exists, because every crossed road \emph{must} point towards 0 for $v$ to reach 0 (the only route from $v$ to 0 in a tree is the tree path).
\begin{lstlisting}[style=cpp]
#include <bits/stdc++.h>
using namespace std;

class Solution {
public:
    int minReorder(int n, vector<vector<int>>& connections) {
        vector<vector<pair<int,int>>> adj(n);        // (neighbour, 1 if the road points away)
        for (auto& c : connections) {
            adj[c[0]].push_back({c[1], 1});
            adj[c[1]].push_back({c[0], 0});
        }
        vector<char> seen(n, 0);
        vector<int> st(1, 0);
        seen[0] = 1;
        int flips = 0;
        while (!st.empty()) {
            int u = st.back(); st.pop_back();
            for (auto& p : adj[u])
                if (!seen[p.first]) {
                    seen[p.first] = 1;
                    flips += p.second;                // road u -> v must be reversed
                    st.push_back(p.first);
                }
        }
        return flips;
    }
};
\end{lstlisting}
\begin{center}
\begin{tikzpicture}[x=1.3cm]
    \node[gvsrc] (c0) at (0,0) {0}; \node[gv] (c1) at (1,0.7) {1}; \node[gv] (c3) at (2,0.7) {3};
    \node[gv] (c2) at (3,0.7) {2}; \node[gv] (c4) at (1,-0.7) {4}; \node[gv] (c5) at (2,-0.7) {5};
    \draw[gpath, -{Stealth[length=3mm]}] (c0) -- (c1); \draw[gpath, -{Stealth[length=3mm]}] (c1) -- (c3);
    \draw[gde] (c2) -- (c3); \draw[gde] (c4) -- (c0); \draw[gpath, -{Stealth[length=3mm]}] (c4) -- (c5);
    \node[note, anchor=west] at (3.5,0) {orange roads point away from 0:\\$0{\to}1$, $1{\to}3$, $4{\to}5$ $\Rightarrow$ \textbf{3}};
\end{tikzpicture}
\end{center}
Samples (checked): $3$, $2$, $0$; a brute force over all $2^{n-1}$ orientations agreed on 3000 random trees. This is the ``search in the right direction'' idea of the theory section turned into a count.
\end{solutionbox}

\begin{solutionbox}[E2.5 Detonate the Maximum Bombs]
\textbf{Model:} vertex $=$ bomb; edge $i \to j$ iff $(x_i - x_j)^2 + (y_i - y_j)^2 \le r_i^2$ --- directed, because the radius belongs to $i$. The answer is $\max_s |R(s)|$: $n$ searches on a graph with at most $n^2 = 10^4$ edges, $\bigO{n^3} = 10^6$.
\begin{lstlisting}[style=cpp]
#include <bits/stdc++.h>
using namespace std;

class Solution {
public:
    int maximumDetonation(vector<vector<int>>& b) {
        int n = b.size();
        vector<vector<int>> adj(n);
        for (int i = 0; i < n; i++)
            for (int j = 0; j < n; j++) {
                if (i == j) continue;
                long long dx = b[i][0] - b[j][0], dy = b[i][1] - b[j][1], r = b[i][2];
                if (dx * dx + dy * dy <= r * r) adj[i].push_back(j);   // 64-bit: up to 2e10
            }
        int best = 0;
        for (int s = 0; s < n; s++) {
            vector<char> seen(n, 0);
            vector<int> st(1, s);
            seen[s] = 1;
            int cnt = 1;
            while (!st.empty()) {
                int u = st.back(); st.pop_back();
                for (int v : adj[u])
                    if (!seen[v]) { seen[v] = 1; cnt++; st.push_back(v); }
            }
            best = max(best, cnt);
        }
        return best;
    }
};
\end{lstlisting}
Sample~3, $[[1,2,3],[2,3,1],[3,4,2],[4,5,3],[5,6,4]]$, builds the edges $0{\to}1, 0{\to}2$; $2{\to}1, 2{\to}3$; $3{\to}1, 3{\to}2, 3{\to}4$; $4{\to}2, 4{\to}3$ (bomb 1 reaches nobody), and $R(0)$ is everything: answer $5$.
\begin{center}
\begin{tikzpicture}[x=1.25cm]
    \node[gvsrc] (d0) at (0,0) {0}; \node[gvdone] (d1) at (1.2,1) {1}; \node[gvdone] (d2) at (1.2,-0.6) {2};
    \node[gvdone] (d3) at (2.6,0.2) {3}; \node[gvdone] (d4) at (4,0.2) {4};
    \draw[gtree, -{Stealth[length=2.6mm]}] (d0) -- (d1); \draw[gtree, -{Stealth[length=2.6mm]}] (d0) -- (d2);
    \draw[gde] (d2) -- (d1); \draw[gtree, -{Stealth[length=2.6mm]}] (d2) to[bend left=15] (d3);
    \draw[gde] (d3) to[bend left=15] (d2); \draw[gde] (d3) -- (d1);
    \draw[gtree, -{Stealth[length=2.6mm]}] (d3) to[bend left=15] (d4); \draw[gde] (d4) to[bend left=15] (d3);
    \draw[gde] (d4) to[bend left=35] (d2);
\end{tikzpicture}
\end{center}
Samples (checked): $2$, $1$, $5$; a closure-based brute force agreed on 3000 random inputs. \textbf{Traps:} building an \emph{undirected} graph (on $[[1,1,5],[5,1,1],[9,1,5]]$ bombs 0 and 2 both reach the small bomb 1 but not each other: the right answer is $2$, while the undirected graph puts all three into one component and claims $3$) and computing $dx^2 + dy^2$ in \texttt{int} ($2\cdot 10^{10}$ overflows).
\end{solutionbox}

\begin{solutionbox}[E2.6 Vertices on some $s$--$t$ walk]
$v$ lies on some walk $s \leadsto v \leadsto t$ iff $v \in R(s)$ and $v \in R^{R}(t)$ (reachable from $t$ in the reverse graph). Two searches, $\VE$:
\begin{lstlisting}[style=cpp]
#include <bits/stdc++.h>
using namespace std;

vector<char> reach(const vector<vector<int>>& g, int s) {
    vector<char> seen(g.size(), 0);
    vector<int> st(1, s);
    seen[s] = 1;
    while (!st.empty()) {
        int u = st.back(); st.pop_back();
        for (int v : g[u]) if (!seen[v]) { seen[v] = 1; st.push_back(v); }
    }
    return seen;
}

int main() {
    int n = 8, s = 5, t = 2;
    int E[][2] = {{0,1},{0,2},{1,3},{2,3},{3,4},{4,1},{5,0},{5,6},{6,7},{7,6}};
    vector<vector<int>> adj(n), rev(n);
    for (auto& e : E) { adj[e[0]].push_back(e[1]); rev[e[1]].push_back(e[0]); }
    vector<char> F = reach(adj, s), B = reach(rev, t);
    int cnt = 0;
    for (int v = 0; v < n; v++)
        if (F[v] && B[v]) { cout << v << " "; cnt++; }
    cout << "-> " << cnt << "\n";                      // 0 2 5 -> 3
}
\end{lstlisting}
Forward from 5 everything is reached; backward from 2 only $\{2, 0, 5\}$. Vertices 1, 3, 4, 6, 7 are reachable from 5 but cannot get back to 2. With $t = 3$ instead, the program finds $\{0,1,2,3,4,5\}$ (6 vertices). A brute force using the transitive closure agreed on 3000 random graphs.
\end{solutionbox}

\begin{solutionbox}[E2.7 Properties Graph]
\textbf{Model:} vertex $=$ array; undirected edge $i - j$ iff $|\mathrm{set}(i) \cap \mathrm{set}(j)| \ge k$. Store every array as a 101-bit \texttt{bitset} of the values it contains --- duplicates collapse by themselves --- so one intersection is \texttt{(has[i] \& has[j]).count()}. \textbf{Counting components with reachability:} scan the vertices; every time a vertex is still unmarked, a new component starts --- count it and mark its whole reach set with one search. Each vertex is marked once, so all searches together cost $\VE$; building the graph costs $\bigO{n^2 \cdot 100/64}$ plus $\bigO{n\,m}$ to read the arrays.
\begin{lstlisting}[style=cpp]
#include <bits/stdc++.h>
using namespace std;

class Solution {
public:
    int numberOfComponents(vector<vector<int>>& properties, int k) {
        int n = properties.size();
        vector<bitset<101>> has(n);                   // has[i][x] = 1 iff x occurs in row i
        for (int i = 0; i < n; i++)
            for (int x : properties[i]) has[i][x] = 1;    // duplicates collapse automatically
        vector<vector<int>> adj(n);
        for (int i = 0; i < n; i++)
            for (int j = i + 1; j < n; j++)
                if ((int)(has[i] & has[j]).count() >= k) {   // distinct common values
                    adj[i].push_back(j);
                    adj[j].push_back(i);
                }
        vector<char> seen(n, 0);
        int components = 0;
        for (int s = 0; s < n; s++) {                 // a new search per unseen vertex
            if (seen[s]) continue;
            components++;
            vector<int> st(1, s);
            seen[s] = 1;
            while (!st.empty()) {
                int u = st.back(); st.pop_back();
                for (int v : adj[u])
                    if (!seen[v]) { seen[v] = 1; st.push_back(v); }
            }
        }
        return components;
    }
};
\end{lstlisting}
Example~1 ($k = 1$): $[1,2]$--$[1,1]$ share 1; $[3,4]$--$[4,5]$--$[5,6]$ form a chain; $[7,7]$ is alone $\Rightarrow$ components $\{0,1\}, \{2,3,4\}, \{5\}$, answer $3$. Example~3 is the trap: $[1,1]$ and $[1,1]$ have one \emph{distinct} common value, less than $k = 2$, so the answer is $2$ (counting matching positions would wrongly join them). The three official examples give $3, 1, 2$ (checked), and a DSU brute force with \texttt{std::set} intersections agreed on 5000 random inputs. Counting components is studied properly in Chapter~4.
\end{solutionbox}
